\section{Introduction}
\label{sec:introduction}

Per-head mixed-precision quantization of the KV cache is a natural strategy for reducing memory during LLM inference: assign fewer bits to insensitive heads and more bits to sensitive ones, achieving lower perplexity than uniform quantization at the same average bit-width.
Recent work on KV cache quantization explores both uniform schemes~\cite{liu2024kivi,hooper2024kvquant,lin2024qserve,he2024zipcache} and token-level mixed-precision approaches~\cite{ge2024fastgen,lei2026arkv}, while complementary methods reduce cache size through eviction~\cite{zhang2023h2o,li2024snapkv}.
RateQuant~\cite{zuo2026ratequant} formalizes per-head allocation as a rate-distortion optimization and shows that greedy allocation is optimal when bit-widths are consecutive integers $\{2,3,4,5,6\}$, exploiting the polymatroid structure of the feasible set (their Proposition~6).

In practice, hardware-efficient inference requires bit-widths drawn from $\{2, 4, 8, 16\}$---the formats natively supported by modern accelerators.
On this non-consecutive set, the polymatroid structure breaks down, and greedy is no longer provably optimal.
One might expect that exact MCKP solvers would fill this gap.
They do not.
On Qwen2.5-7B at $4\times$ compression, unconstrained greedy achieves 32.45 PPL and unconstrained MCKP achieves 36.92---both \emph{far worse} than uniform 4-bit (15.45).
The optimization does not merely underperform; it fails catastrophically.

The root cause is not solver suboptimality but a violation of the objective's additive distortion assumption.
Both solvers assign 68--75 of 112 KV heads to 2-bit, where per-head quantization errors are large enough to compound through the residual stream.
The additive model predicts $\sim$11 PPL for these allocations; the actual PPL is $3$--$4\times$ higher.
At 4-bit and above, per-head distortions are small and compound negligibly---uniform 4-bit achieves close to its additive prediction---but 2-bit distortions are orders of magnitude larger.
When most heads are quantized to 2-bit simultaneously, cumulative truncation errors exceed the network's capacity to compensate, and perplexity degrades super-linearly.

This observation points to a minimal intervention: cap the number of 2-bit assignments.
We add a single cardinality constraint $\sum_h \mathbb{1}[b_h{=}2] \leq k$ to the standard MCKP formulation (\S\ref{sec:method}).
On Qwen2.5-7B, setting $k{=}8$ (7\% of heads) yields 11.23 PPL---a 27.3\% reduction over uniform 4-bit.
The constraint works by forcing the solver to use 2-bit only on the least-sensitive heads, while redistributing the freed budget to protect sensitive heads at 8-bit or 16-bit.
This is not a heuristic patch: the constraint directly addresses the identified failure mode and is compatible with any solver.

A surprising consequence is that, under the cardinality constraint, the choice of solver becomes irrelevant.
Constrained greedy and constrained MCKP produce identical allocations---not approximately, but exactly: all 112 heads receive the same bit-width on Qwen2.5-7B, all 256 on Llama-3.1-8B, and 255 of 256 on Mistral-7B-v0.3 (the single difference is a zero-distortion tie with no PPL impact).
The constraint narrows the feasible set enough that greedy's locally optimal decisions are globally optimal (\S\ref{sec:solver_equiv}), enabling practitioners to use the simpler $O(H \log H)$ greedy algorithm without sacrificing solution quality.

However, non-uniform allocation does not help universally.
On Llama-3.1-8B, constrained MCKP achieves 8.45 PPL vs.\ uniform 8.47---a marginal 0.2\% improvement.
On Mistral-7B-v0.3, it is slightly \emph{worse}: 7.16 vs.\ 7.13.
The difference across architectures traces to head sensitivity heterogeneity: Qwen's KV heads vary widely in 2-bit distortion, with Layer-0 heads contributing a disproportionate share, while Mistral's heads are nearly uniform in sensitivity.
When there is no heterogeneity to exploit, non-uniform allocation adds noise without benefit.

This motivates a diagnostic question that prior work does not address: \emph{when} should a practitioner use per-head mixed-precision quantization?
We propose a heterogeneity diagnostic based on the top-5\% concentration of per-head distortion (\S\ref{sec:heterogeneity}).
On our three architectures, the metric cleanly separates the beneficial case (Qwen, concentration 0.59--0.66, 27.3\% PPL improvement) from the neutral case (Llama, 0.34--0.59, marginal) and the counterproductive case (Mistral, 0.42, slightly worse than uniform).
The diagnostic reuses profiling data already required by the solver and adds no computational overhead.

Our contributions are:
\begin{enumerate}[nosep]
\item We identify 2-bit distortion compounding as the cause of catastrophic failure in unconstrained per-head KV cache quantization on hardware-aligned bit-widths, and show that a cardinality constraint on 2-bit assignments is a minimal sufficient fix, reducing Qwen2.5-7B PPL from 32--37 (unconstrained) to 11.23 (constrained) at $4\times$ compression (\S\ref{sec:compounding}, \S\ref{sec:experiments}).

\item We demonstrate that the cardinality constraint makes solver choice irrelevant: constrained greedy and constrained MCKP produce identical per-head allocations across three architectures, enabling use of a simple greedy algorithm (\S\ref{sec:solver_equiv}).

\item We propose a heterogeneity-based diagnostic framework that predicts whether non-uniform allocation will help a given architecture, validated on three models spanning beneficial, marginal, and counterproductive outcomes (\S\ref{sec:heterogeneity}).
\end{enumerate}
